\section{Additional Exercises}

\subsection{Practical Rules For How $\nabla$ Acts}

\begin{exercise}\label{ex:connections-1}
What is the result of the following applications of a torsion free covariant derivative?
\begin{itemize}
    \item There are some others in the video, but they are basically covered in the lectures so I won't repeat them here.
    \item $\big(\nabla_{[m} A\big)_{n]}$,
    \item $\big(\nabla_{[m}\omega\big)_{nr]}$.
\end{itemize}
\end{exercise}
\hyperref[sol:connections-1]{\small \textit{Solution on p.~\pageref*{sol:connections-1}}}

\subsection{Connection Coefficients}

\textit{There is a question about stating the chart transformation laws of the $\Gamma$s and stating which class of transformations make the $\Gamma$s look like tensors. We have discussed this in the notes already, but I have put this here to remind readers to re-read \Cref{rem:GammasTensorTransformation} as it's an important point often missed.}

\begin{exercise}\label{ex:connections-2}
Let $(\cM,\cO,\cA,\nabla)$ be the flat plane. Consider two charts that both cover the upper half plane, that is all the points $(a,b)\in\R^2$ with $b>0$, one representing the familiar Cartesian coordinates and the other the familiar polar coordinates on there.

We already know the chart transition map from Cartesian to polar coordinates is given by
\bse
    y\circ x^{-1}(a,b) = \bigg( \sqrt{a^2+b^2}, \arccos\bigg(\frac{a}{\sqrt{a^2+b^2}}\bigg)\bigg),
\ese
while the inverse transition map from polar to Cartesian is given by
\bse
    x\circ y^{-1}(r,\varphi) = \big( r\cos\varphi, r\sin\varphi), \qquad \text{for } r\in\R^+ \text{ and } \varphi\in(0,\pi).
\ese
Starting from the assumption of vanishing connection coefficient functions in the Cartesian chart, calculate the connection coefficient functions $\Gamma^a_{(y)bc}$ w.r.t. the polar chart!
\end{exercise}
\hyperref[sol:connections-2]{\small \textit{Solution on p.~\pageref*{sol:connections-2}}}
